which shows that @@M0@@.

\sgaStableTarget{sga6.e05.cor.6-12}\textbf{Corollary 6.12.} \uline{If} @@M1@@ \uline{is of characteristic zero}, \uline{and} @@M2@@ \uline{is a} @@M3@@-@@M4@@-\uline{algebra} \uline{with discrete filtration}, \uline{then the completed Chern class}:
@@M5@@
\uline{is an isomorphism}.

The proof is contained in that of 6.11.

\sgaStableTarget{sga6.e05.def.6-13}\textbf{Definition 6.13.} \uline{Let} @@M6@@ \uline{be a graded} @@M7@@-\uline{algebra}, \uline{and} @@M8@@ \uline{an} \uline{augmented} @@M9@@-@@M10@@-\uline{algebra}; \uline{we call a Chern theory for} @@M11@@ \uline{with values in} @@M12@@ \uline{any} @@M13@@-\uline{homomorphism}: @@M14@@, \uline{compatible with the augmentations}.

We define a functor on the category of graded @@M15@@-algebras, with values in the category of sets, by associating with every graded @@M16@@-algebra @@M17@@ the set of Chern theories for @@M18@@ with values in @@M19@@, namely @@M20@@.

\sgaStableTarget{sga6.e05.prop.6-14}\textbf{Proposition 6.14.} \uline{The functor defined in} 6.13 \uline{is representable}. \uline{If} @@M21@@ \uline{is} \uline{of characteristic zero}, \uline{and} @@M22@@ \uline{has discrete filtration}, \uline{it is represented} \uline{by the graded ring associated with} @@M23@@.

Let us first prove the first assertion. We consider the graded @@M24@@-algebra generated by generators @@M25@@, with @@M26@@, @@M27@@, where @@M28@@ has degree @@M29@@, and take its quotient by the following relations:
@@M30@@

